\scanpage{60}


\mspage{3}

L'énoncé (i bis) est dual de (i).

\uline{Corollaire 1} \inserted{Sous les cond. de prop. 1.} Supposons que $i_!$ existe (condition (i) de la proposition 1). Les conditions suivantes sont équivalentes

a) $j^*$ exact, $i_!\simeq\operatorname{Ker}j^*$ [à fortiori, $i_!$ \uline{est exact}\deleted{\unsure{et pl. fid}}]

b) $i_!$ a une image \inserted{essentielle} épaisse [épaisse à g. suffit]

\deleted{\unsure{b}}c) $i_!i^*(F)\unsure{\to|-}F$ est injectif pour tt $F$ i.e. la suite
\[
0\longrightarrow\underset{\substack{\shortparallel\\F^U}}{i_!i^*F}
\longrightarrow F\longrightarrow\underset{\substack{\shortparallel\\F^Y}}{j_*j^*F}
\longrightarrow 0\qquad\text{est exacte}
\]

d) $j^*$ \uline{est exact}

\deleted{Alors $i_!\simeq\operatorname{Ker}j^*$, à fortiori $i_!$ est exact}

a) $\Rightarrow$ b) trivial,

b) $\Rightarrow$ c)
\[
\left\{\begin{aligned}
&\text{Car }0\to N\to i_!i^*F\to F \overset{i^*\text{ exact}}{\Rightarrow}0\to i^*N\to\underbrace{i^*i_!}\,\underbrace{i^*F}\to\underbrace{i^*F}\quad\substack{\text{dernière fl. un isom.}\\\text{($i^*$ passage au quotient)}}\\
&\text{donc }i^*(N)=0\ \deleted{\unsure{\text{donc}}},\ \text{or } N=i_!(M)\ \text{par b)},\quad i^*i_!(M)=0\Rightarrow M=0,\Rightarrow N=0
\end{aligned}\right.
\]

c) $\Rightarrow$ a) \deleted{\unsure{$F\to F'\to F''\to 0$}} Supposons \uline{$0\to F\to F'\to F''\to0$} d'où

\deleted{a) $\Rightarrow$ d) trivial}\quad\deleted{\unsure{d) $\Rightarrow$ a)|d) $\Rightarrow$ b)}}

\diagram{\begin{tikzcd}[ampersand replacement=\&,column sep=small,row sep=small]
{} \& {} \& \arrow[d] {} \& {} \& {}\\
\deleted{0}\arrow[r] \& i_!i^*F\arrow[r]\arrow[d] \& F\arrow[r]\arrow[d] \& j_*j^*F\arrow[r]\arrow[d] \& 0\\
\deleted{\unsure{0\to}} \& i_!i^*F'\arrow[r]\arrow[d] \& F'\arrow[r]\arrow[d] \& j_*j^*F'\arrow[r] \& 0\\
0\arrow[r] \& i_!i^*F''\arrow[r] \& F''
\end{tikzcd}}

et utilisant l'injectivité de $i_!i^*F''\to F''$, on trouve $j_*j^*F\to j_*j^*F'$ est injectif, donc $j^*F\to j^*F'$ inj. d'où $j^*$ ex. (car il est déjà exact à droite).

Soit $F=i_!G$, alors $i_!i^*F\to F$ est un isom [$i^*$ passage au quotient] donc $j_*j^*F=0$ \deleted{\unsure{est un isom}} par la suite exacte (i)\inserted{ déjà vu dans prop. 1.} donc $j^*F=0$. Réciproquement, si $j^*F\unsure{=0|-0}$, alors par c), on trouve $i_!i^*F\xrightarrow{\sim}F$, i.e. $F\unsure{\in|\equiv}\operatorname{Im}\,\mathrm{ess}\,i_!$

% archive sheet number: 59
